% electric_block.tex -- the derivation section.
%
% Load-bearing equations are \input from sections/generated/*.tex, which are
% emitted by src/utilities/electric_induced_action.py via sympy.latex() -- so
% the displayed equations are provably identical to what the verifier computed
% (CLAUDE.md "SymPy-generated equations"). Do not hand-edit those fragments.
%
% Backs the electric-block, covariant-completion, and birefringence-verdict rows
% of the audit table (Table~\ref{tab:audit}); verified by
% src/utilities/electric_induced_action.py and src/experiments/exp_01_*.py.

\label{sec:electric_block}

The tick rule couples the two static fields on different footings, and that
asymmetry is the whole reason this block is new theory. A magnetic field enters
as a \emph{spatial} link phase: the hop from a site to a neighbour along a
lattice direction $V_a$ picks up the Peierls phase $\exp(i\,A\!\cdot\!V_a)$, so a
closed spatial plaquette carries a gauge-invariant Wilson-loop holonomy. An
electric field enters instead as an \emph{on-site} phase advance
$\delta\phi(x)=\omega+V(x)$---the temporal potential $A_0$ added to the rest
phase $\omega$---with no spatial link and hence, at first sight, no loop to read
a holonomy from. Consequently the electric action block, the permittivity
$\varepsilon$, appears in neither Paper~I nor Paper~II~\cite{PaperI,PaperII}, and
because the framework establishes only the spatial point group $O_h$ and not the
Lorentz boosts, there is no symmetry route from the magnetic $\mu^{-1}$ to the
electric $\varepsilon$. We derive it directly, then show the two blocks are bound
by an exact algebraic identity that decides the birefringence question.

\subsection{The magnetic anchor}
\label{subsec:magnetic_anchor}

We first fix the object the electric block must mirror. A spatial plaquette
spanned by two hop directions $V_a,V_b$ has, for a uniform field, the holonomy
$(V_a\!\times\!V_b)\!\cdot\!B$; summing the squared holonomy over the three
independent bipartite plaquettes gives the induced $O(B^2)$ action
$\text{density}_B=\sum_{a<b}\bigl((V_a\!\times\!V_b)\!\cdot\!B\bigr)^2=B^{\!\top}
Q_B\,B$, with
%
\begin{equation}
    \label{eq:magnetic_Q}
    \input{sections/generated/magnetic_Q}
\end{equation}
%
whose eigenvalues are $\{4,4,16\}$: the optical axis $(1,1,-1)$ is
\emph{enhanced} (eigenvalue $16$) and the perpendicular plane is doubly
degenerate at $4$. This reproduces Paper~I, Appendix~B, and is reproduced
numerically by \texttt{dcl-core}'s \texttt{exp\_04} to machine
precision~\cite{PaperI,dclcore}; it is retained throughout as an exact
consistency anchor (Table~\ref{tab:audit}).

\subsection{The electric block}
\label{subsec:electric_block}

The pessimistic reading---``no electric Wilson loop''---is too strong. The
on-site advance $\delta\phi=\omega+V(x)$ \emph{is} a temporal link: it joins a
site to itself at the next tick. A \emph{temporal} plaquette spanned by one hop
direction $V_a$ and the tick direction therefore closes on the bipartite lattice,
and for a uniform static field ($V(x)=-E\!\cdot\!x$) its holonomy is the on-site
phase difference along the hop,
%
\begin{equation}
    \label{eq:temporal_holonomy}
    \Theta_a(x) = \delta\phi(x) - \delta\phi(x+V_a)
                = V(x) - V(x+V_a) = (V_a\!\cdot\!E)\,a_t,
\end{equation}
%
where the rest phase $\omega$ cancels around the loop and $a_t$ is the tick's
temporal extent. Where the magnetic loop needs \emph{two} spatial directions
(and so a cross product), the electric loop needs only \emph{one}. Summing the
squared holonomy over the three hop directions gives the induced $O(E^2)$ action
$\text{density}_E=\sum_a\Theta_a^2 = a_t^2\,E^{\!\top}\varepsilon\,E$, with the
permittivity block
%
\begin{equation}
    \label{eq:electric_P}
    \input{sections/generated/electric_P}
\end{equation}
%
Its eigenvalues are $\{1,4,4\}$: the same optical axis $(1,1,-1)$ is now
\emph{suppressed} (eigenvalue $1$), with the perpendicular plane again degenerate
at $4$. The electric block is thus the exact mirror of the magnetic one---axis
suppressed where the magnetic axis is enhanced---and its sense reproduces the
axis-suppression sign already seen in Paper~IV's static screen
\texttt{exp\_03a}~\cite{PaperIV}.

This block is not merely posited from the continuum holonomy: it is read off the
engine's actual tick evolution. Seeding a uniform, real probe state and taking one
step of the hop operator, the on-site term $i\sin(\delta\phi/2)\,\psi_R$ is purely
imaginary while the kinetic hop is purely real, so
$\operatorname{Im}(\psi_R^{\text{new}})/\psi_R=\sin(\delta\phi/2)$ recovers
$\delta\phi=\omega+V(x)$ directly from the step output (to $10^{-19}$); the
permittivity $\varepsilon$ of Eq.~\eqref{eq:electric_P} is then extracted from
that recovered phase (\texttt{exp\_01}, the temporal-sector mirror of
\texttt{exp\_04}). A dropped potential collapses $\varepsilon$ to zero and a sign
or tick-weight error trips the fidelity check, so the result genuinely depends on
the engine's coupling; the recovered unit tick weight ($a_t=1$) is a measured
fact, not an assumption.

\subsection{The covariant completion}
\label{subsec:covariant_completion}

The two blocks are not independent. Identifying the macroscopic response tensors
as the permittivity $\varepsilon=P$ and the inverse permeability $\mu^{-1}=Q_B$,
they are bound by an exact algebraic identity.

\begin{theorem}[Adjugate relation]
    \label{thm:adjugate}
    For any three spatial hop vectors, the magnetic induced-action tensor is the
    adjugate of the electric one:
    %
    \begin{equation}
        \label{eq:adjugate}
        \input{sections/generated/adjugate_relation}
    \end{equation}
\end{theorem}

\begin{proof}
    Arrange the hop vectors as the columns of a $3\times3$ matrix $M$. Then
    $\varepsilon=\sum_a V_a V_a^{\!\top}=M M^{\!\top}$, while the cross products
    $V_a\!\times\!V_b$ are the rows of $\operatorname{adj}(M)$ (the columns of
    $\operatorname{adj}(M)^{\!\top}$), so
    $Q_B=\sum_{a<b}(V_a\!\times\!V_b)(V_a\!\times\!V_b)^{\!\top}
    =\operatorname{adj}(M)^{\!\top}\!\operatorname{adj}(M)
    =\operatorname{adj}(M^{\!\top})\operatorname{adj}(M)
    =\operatorname{adj}(MM^{\!\top})=\operatorname{adj}(\varepsilon)$. The
    identity is verified symbolically for general vector components in
    \texttt{src/utilities/electric\_induced\_action.py}.
\end{proof}

\begin{corollary}[Isotropic impedance product]
    \label{cor:impedance}
    Because the adjugate carries each eigenvalue of $\varepsilon$ to the product
    of the other two, every principal direction $i$ obeys
    $\varepsilon_i\,(\mu^{-1})_i=\varepsilon_i\,(\operatorname{adj}\varepsilon)_i
    =\det\varepsilon$. The impedance product is therefore isotropic even though
    $\varepsilon$ and $\mu^{-1}$ are each anisotropic, with opposite senses about
    the shared axis $(1,1,-1)$.
\end{corollary}

For the lattice values this reads $\varepsilon_a\,\mu^{-1}_a=1\cdot16=16$ on the
axis and $\varepsilon_p\,\mu^{-1}_p=4\cdot4=16$ in the perpendicular plane: the
electric axis-suppression (a factor $\tfrac14$) exactly reciprocates the magnetic
axis-enhancement (a factor $4$).

\subsection{The birefringence verdict}
\label{subsec:birefringence}

Corollary~\ref{cor:impedance} is exactly the condition that removes vacuum
birefringence. For a plane wave with wavevector $k$ in an anisotropic medium with
permittivity $\varepsilon$ and inverse permeability $\mu^{-1}$, Maxwell's
equations give the dispersion relation $\det(K\,\mu^{-1}K+\omega^2\varepsilon)=0$,
where $K$ is the curl operator $[k]_\times$. With $\mu^{-1}=\operatorname{adj}
(\varepsilon)$ this determinant factorizes---for a \emph{general} symmetric
$\varepsilon$ and general $k$, requiring no shared eigenvectors---as
%
\begin{equation}
    \label{eq:dispersion}
    \input{sections/generated/dispersion_factorization}
\end{equation}
%
The $\omega^2=0$ root is the unphysical longitudinal mode; the physical root is
the \emph{doubled} transverse root, so both photon polarizations share the single
dispersion
%
\begin{equation}
    \label{eq:shared_dispersion}
    \input{sections/generated/shared_dispersion}
\end{equation}
%
and the polarization splitting vanishes identically,
%
\begin{equation}
    \label{eq:split}
    \input{sections/generated/birefringence_condition}
\end{equation}
%
Gauge-sector vacuum birefringence \emph{cancels}. Because the splitting depends
only on the anisotropy ratios---fixed by the geometry through
Theorem~\ref{thm:adjugate}---and not on the overall block scales, the cancellation
is invariant under independent rescaling of $\varepsilon$ and $\mu^{-1}$: it is
immune to the undetermined tick-weight and $1/g^2$ coupling factors, and it holds
for any hop vectors. A numerical sweep over $2\times10^4$ random propagation
directions confirms the vanishing split to $\sim10^{-15}$, guarding the symbolic
proof. This is the verdict Paper~IV waits on: the framework predicts the
\emph{absence} of gauge-sector vacuum birefringence about $(1,1,-1)$.

\subsection{The optical axis, the speed anisotropy, and \texorpdfstring{$O_h$}{Oh}
restoration}
\label{subsec:scope_verdict}

Two features deserve statement rather than concealment. The first is geometric and
explains where the optical axis comes from. As \emph{axes}, the three hop directions
$V_1,V_2,V_3$ are three of the four cube body-diagonals; the optical axis
$(1,1,-1)$ is precisely the \emph{fourth} diagonal, the one the hop set omits.
Selecting three of the four diagonals breaks the cubic group $O_h$ down to the
trigonal group $D_{3d}$ about the omitted diagonal, and a $D_{3d}$-invariant
quadratic form is exactly uniaxial---which is why both blocks come out
$\{\text{axis};4,4\}$. The axis is not imposed; it is the diagonal the lattice leaves
out.

The second is the shared dispersion $\omega^2=k^{\!\top}\varepsilon\,k$ of
Eq.~\eqref{eq:shared_dispersion}, which is direction-dependent: the common phase
speed $v^2(\hat k)$ ranges over the eigenvalues of $\varepsilon$, a factor of about
two between the axis and the perpendicular plane. This is a \emph{common-mode} speed
anisotropy shared by both polarizations---not birefringence---but it is not small.
Its status is clarified by the same geometry. There are four diagonal domains, one
trigonal about each body-diagonal; averaging the two blocks over them restores full
cubic symmetry exactly,
%
\begin{equation}
    \label{eq:oh_average}
    % AUTO-GENERATED by src/utilities/electric_induced_action.py -- do not edit.
\langle\varepsilon\rangle = \left[\begin{matrix}3 & 0 & 0\\0 & 3 & 0\\0 & 0 & 3\end{matrix}\right],\qquad \langle\mu^{-1}\rangle = \left[\begin{matrix}8 & 0 & 0\\0 & 8 & 0\\0 & 0 & 8\end{matrix}\right]

\end{equation}
%
so the $O_h$-restored medium is isotropic Maxwell, with no speed anisotropy. The
decisive observation is that the adjugate relation---and hence the birefringence
cancellation---holds \emph{in every single domain}, non-trivially, and equally after
averaging, trivially. The cancellation is therefore robust to whether the physical
vacuum is a single trigonal domain or an $O_h$-restored average; the speed anisotropy
is the one feature that distinguishes the two, and however that question resolves it
cannot revive birefringence. Deciding it---single-domain, with the factor-of-two
speed anisotropy as a genuine (and observationally constrained) prediction, versus
$O_h$-restored and isotropic---is a question about vacuum structure separate from,
and not gating, the polarization verdict. (A remaining bookkeeping item, immaterial
to all of the above, is the relative normalization of $\varepsilon$ and $\mu^{-1}$:
the tick-weight and one-loop $1/g^2$ factors set the overall photon speed but, being
an independent rescaling of each block, cannot affect the birefringence verdict.)
Accordingly the electric block and its covariant completion are recorded as complete,
while the unconditional birefringence verdict is held partial pending this
vacuum-structure question and Paper~IV's full large-$N$ dispersion classification
(Table~\ref{tab:audit}).
